\documentclass[11pt,a4paper]{ctexart}

% Page geometry
\usepackage[top=2.5cm,bottom=2.5cm,left=2.8cm,right=2.8cm]{geometry}

% Minimal typography
\usepackage{setspace}
\onehalfspacing
\usepackage{enumitem}
\setlist{nosep,leftmargin=*}

% Tables
\usepackage{booktabs}
\usepackage{array}

% Colors (minimal)
\usepackage{xcolor}
\definecolor{accent}{RGB}{45,55,72}
\definecolor{muted}{RGB}{120,130,140}

% Section styling
\usepackage{titlesec}
\titleformat{\section}{\Large\bfseries\color{accent}}{\thesection}{1em}{}
\titleformat{\subsection}{\large\bfseries\color{accent}}{\thesubsection}{1em}{}
\titleformat{\subsubsection}{\normalsize\bfseries\color{accent}}{\thesubsubsection}{1em}{}
\titlespacing*{\section}{0pt}{1.8em}{0.8em}
\titlespacing*{\subsection}{0pt}{1.2em}{0.5em}

% Links
\usepackage[colorlinks=true,linkcolor=accent,urlcolor=muted]{hyperref}

% Code / verbatim
\usepackage{listings}
\lstset{
  basicstyle=\small\ttfamily,
  breaklines=true,
  frame=none,
  backgroundcolor=\color{gray!8},
  xleftmargin=1em,
  xrightmargin=1em,
  aboveskip=0.8em,
  belowskip=0.8em,
}

% Header/Footer (minimal)
\usepackage{fancyhdr}
\pagestyle{fancy}
\fancyhf{}
\fancyfoot[C]{\small\color{muted}\thepage}
\renewcommand{\headrulewidth}{0pt}

\begin{document}

% --------------------------------------------------
% 标题页
% --------------------------------------------------
\begin{center}
\vspace*{1em}
{\LARGE\bfseries QQClawNewbie}\\[0.6em]
{\Large 个性化 OpenClaw 配置系统（Demo）}\\[1.5em]
\end{center}

\noindent
QQClawNewbie 是一款面向移动端 H5 的个性化 AI 助手配置系统 Demo。它通过多通道、可视化、可交互的配置流程，帮助 QQ 生态内的用户以低门槛、高趣味性的方式完成 OpenClaw 人格配置——不需要阅读任何技术文档，点选标签、聊天对话、导入文件或授权数据，即可快速生成一份带有人格类型、标签云、Soul 设定与原始 JSON 的完整配置。

当前项目是一个可运行的前端 Demo：打开页面，选择任意配置通道，完成交互后即可在结果页查看并复制生成的配置。

\vspace{1em}

% --------------------------------------------------
\section{本 Demo 的功能}
% --------------------------------------------------

\subsection{面向人群}

\begin{itemize}
  \item QQ 生态内的普通用户（对 OpenClaw 感到陌生，希望快速上手）。
  \item 有明确需求的进阶用户（知道自己想要什么类型的助手，希望快速模板化配置）。
  \item 已有配置或聊天记录的迁移用户（希望把历史偏好一键导入新助手）。
  \item 愿意基于自身数据做深度优化的用户（希望通过授权数据获得更精准的推荐）。
\end{itemize}

\subsection{解决的痛点}

\begin{itemize}
  \item \textbf{陌生感}：用户不知道 OpenClaw 能做什么，也不知道从何开始配置。
  \item \textbf{配置门槛高}：传统表单式配置需要理解人格特质、能力插件、Soul 设定等概念，认知成本大。
  \item \textbf{缺乏趣味性}：纯文本表单的填写过程枯燥，用户容易中途流失。
  \item \textbf{迁移困难}：用户可能在其他平台已有偏好记录，但无法快速迁移到 OpenClaw。
\end{itemize}

\subsection{已实现的核心功能}

\begin{enumerate}[label=\arabic*.]
  \item \textbf{配置大厅首页}：顶部渐变品牌区 + 四通道卡片展示；实时进度追踪（Zustand 全局状态）；一键生成与重置。
  \item \textbf{通道 A — 角色模板}：24 个标签覆盖使用场景、性格倾向、交互偏好、关注领域四大维度；多选高亮 + 数量统计 + 一键清空；确认后自动返回大厅。
  \item \textbf{通道 B — 对话定制}：拟真 IM 界面（头像、在线状态、打字指示器）；5 轮固定对话流程覆盖场景、风格、创意、回应等维度；快捷回复 + 手动输入；自动提取标签存入状态。
  \item \textbf{通道 C — 文件导入}：拖拽/点击上传区域；模拟上传 → 解析 → 成功流程；按文件类型自动生成标签；隐私安全提示（仅本地解析）。
  \item \textbf{通道 D — 数据授权}：完整 5 步授权流程（介绍 → 权限选择 → 扫描分析 → 分析报告 → 完成）；雷达扫描动画 + 6 维分析报告；用户可自主选择采纳标签。
  \item \textbf{结果页}：展示匹配人格类型与副标题；配置完成度可视化；动态标签云；OpenClaw 4 大维度概览；Soul 设定（Markdown，可复制）；原始 JSON（可复制）；重新配置入口。
  \item \textbf{预设人格库}：基于对 OpenClaw 官方 best-practices（\url{https://openclaw-ai.com/en/blog/best-practices}）的调研，已沉淀 9 个高质量预设（极客架构师、灵感捕手、生活管家、战略军师、学术研究员、代码导师、数据分析师、产品经理、心灵疗愈师），覆盖九大典型场景。
  \item \textbf{标签匹配算法}：基于用户在各通道产生的标签，与 9 类预设标签做交集计数匹配，返回最接近的人格；无标签时返回通用配置。
\end{enumerate}

% --------------------------------------------------
\section{目标版本 vs 当前 Demo 的差异说明与后续计划}
% --------------------------------------------------

\subsection{当前 Demo 已经能做什么}

\begin{itemize}
  \item 配置大厅展示四通道卡片，实时追踪各通道完成状态，支持一键重置。
  \item 通道 A 提供 4 大维度共 24 个标签的多选界面，标签数据汇入全局状态。
  \item 通道 B 模拟 5 轮拟真对话，用户通过快捷按钮或自由输入完成交互，系统自动提取标签。
  \item 通道 C 模拟拖拽上传、解析流程与标签提取，展示隐私安全提示。
  \item 通道 D 完成完整的 5 步授权引导流程（介绍→权限→扫描→报告→完成），含雷达扫描动画。
  \item 结果页展示匹配人格类型、配置完成度、动态标签云、4 大维度概览、Soul Markdown 与原始 JSON。
  \item 预设人格库已沉淀 9 个高质量预设，覆盖技术、创意、生活、战略、学术、教学、数据、产品、心理九大场景。
  \item 基于标签的静态匹配算法已可运行，根据用户选择返回最接近的预设人格。
\end{itemize}

\subsection{与目标版本的主要差距}

\begin{itemize}
  \item \textbf{缺少真实大模型对话}：通道 B 的 5 轮对话为固定文案，目标版本将接入 GPT/Claude 等 API，实现真正的多轮理解与动态生成。
  \item \textbf{缺少真实文件解析}：通道 C 为随机模拟结果，目标版本将使用浏览器 File API 读取本地文件，完成 JSON/YAML/Markdown 的结构化解析。
  \item \textbf{缺少真实数据授权}：通道 D 完全模拟，目标版本在合规授权下可接入真实行为数据，生成个性化报告。
  \item \textbf{配置生成过于静态}：当前为标签匹配预设模板，目标版本将接入大模型进行深度理解，综合多通道输入生成真正个性化的配置。
  \item \textbf{缺少配置导出与分享}：当前仅支持页面内复制 JSON，目标版本将支持导出文件、生成分享链接。
  \item \textbf{缺少试聊与预览}：结果页无即时交互，目标版本将增加“试聊”入口，让用户即时检验配置效果。
  \item \textbf{缺少多端适配}：当前为纯 H5，目标版本将扩展至微信小程序或桌面端。
\end{itemize}

\subsection{后续优化目标与实现计划}

\subsubsection{1) 短期改进（1–2 周内可完成）}

\begin{itemize}
  \item \textbf{扩充预设人格库}：至少再增加 4–6 个高质量预设，覆盖更多垂直场景。
  \item \textbf{优化匹配算法}：从“包含即匹配”升级为加权匹配算法，考虑标签权重与通道优先级。
  \item \textbf{丰富对话分支}：即使未接入大模型，也可通过更丰富的条件分支让对话体验更自然。
  \item \textbf{增加导出功能}：支持 JSON 文件下载与分享链接生成。
\end{itemize}

\subsubsection{2) 中期改进（1 个月内可完成）}

\begin{itemize}
  \item \textbf{接入真实大模型对话}：将通道 B 替换为真正的 AI 多轮对话，这是体验提升最大的改进点。
  \item \textbf{实现真实文件解析}：使用浏览器 File API 读取并解析用户本地文件。
  \item \textbf{增加实时预览}：结果页增加“试聊”入口，即时检验配置效果。
  \item \textbf{UI 细节打磨}：深色模式、骨架屏、空状态引导、加载优化（图片压缩、懒加载、代码分割）。
\end{itemize}

\subsubsection{3) 长期愿景（课程结束后可继续探索）}

\begin{itemize}
  \item \textbf{构建配置市场}：用户可分享或下载他人的 OpenClaw 配置模板，形成社区生态。
  \item \textbf{多端适配}：从纯 H5 扩展至微信小程序或桌面端。
  \item \textbf{接入真实数据授权}：在合规允许下基于真实使用数据实现精准推荐。
  \item \textbf{工程化完善}：接入错误监控、性能分析、单元测试与 E2E 测试覆盖。
\end{itemize}

\end{document}
